% !PRINTED 262
\weberpage{262}

%<<CONT
\weq{1}{\psi = \varphi(x_1, x_2, x_3 \ldots x_m) - \lambda x_m^2}
vanishes. For the determinant is
\[
\begin{vmatrix}
a_{1,1},\, a_{1,2} \ldots a_{1,m} \\
a_{2,1},\, a_{2,2} \ldots a_{2,m} \\
\cdots\cdots\cdots\cdots \\
a_{m,1},\, a_{m,2} \ldots a_{m,m} - \lambda
\end{vmatrix}
= R_m - \lambda R_{m-1},
\]
and it vanishes when we set
\[
\lambda = \frac{R_m}{R_{m-1}}
\]

The function $\psi$ can then be expressed in terms of $m-1$ variables $y$, and by formula (6), \S.\,81, one obtains
\[
\psi(x_1, x_2 \ldots x_m) = \psi(y_1, y_2 \ldots y_{m-1}, 0),
\]
or, by (1),
\weq{2}{\varphi(x_1, x_2 \ldots x_m) = \frac{R_m}{R_{m-1}}\, x_m^2 + \varphi(y_1, y_2 \ldots y_{m-1}, 0),}
and the investigation of the function $\varphi$ of $m$ variables is thereby reduced to the investigation of $\varphi_1(y) = \varphi(y_1, y_2 \ldots y_{m-1}, 0)$ of $m-1$ variables, whose determinant is equal to $R_{m-1}$ and hence different from zero. According as $R_m : R_{m-1}$ is positive or negative, the function $\varphi(x)$ will have one more positive or one more negative square than $\varphi_1(y)$.

By applying the same procedure to $\varphi_1(y)$ and the following functions, we obtain the theorem:

The number of positive and negative terms of the sequence
\weq{3}{\frac{R_m}{R_{m-1}}, \quad \frac{R_{m-1}}{R_{m-2}}, \ldots \frac{R_1}{R_0}}
agrees with the number of positive and negative squares into which the function $\varphi(x)$ can be decomposed.

We wish to give this theorem a somewhat different expression, but first premise the following explanation.

If a sequence of real numbers different from zero is given in a definite arrangement, the signs of these quantities can change in manifold ways; if two quantities of the same sign follow one another, there is a \gesperrt{succession of signs} (permanence), but if a quantity is followed by another of opposite sign, we have a \gesperrt{change of sign} (variation).
